
% Authoritative LaTeX manuscript; Markdown is a generated compatibility projection.

Execution and placement require two related views: where components run and connect, and how a host schedules their work. Keeping those views separate lets review distinguish changes to topology from changes to execution policy.

\textbf{Engineering question.} Where does work run, and which scheduling decisions may change with the host it runs on?

\textbf{Representation.} Use separate views for topology and host policy.

\begin{itemize}
\tightlist
\item
  \textbf{Deployment topology} represents where components run and how they connect: web → queue → worker → storage.
\item
  \textbf{Host execution policy} maps each host profile to a scheduling policy: elastic fan-out, bounded fan-out, or serialized execution.
\end{itemize}

\autoref{fig-g4-execution-placement} shows the stable graph with the policy as a separate layer beneath it.

\begin{figure}[htbp]
\centering
\includegraphics[width=\linewidth]{book/latex/assets/appendix-g-4-host-execution-policy-c72a07bdc9.pdf}
\caption{\emph{Stable topology, variable execution policy.} Deployment topology remains stable while host profiles select different scheduling behavior.}
\label{fig-g4-execution-placement}
\end{figure}

\textbf{Property.} A host\textquotesingle s concurrency ceiling belongs to scheduling policy, not to the deployment graph.

\textbf{Enforcement and correspondence.} Deployment configuration reconciles against the topology; scheduler behavior reconciles against the host execution policy. Either finding may remain advisory or feed a deployment or scheduler gate.
